\chapter{System Design}\label{chap:design}

\section{Design Goals and Principles}

The design of \ProjectName{} follows from the requirements of \cref{chap:requirements} and the principles reviewed in
\cref{chap:literature}. Seven principles guided every decision:

\begin{itemize}
  \item \textbf{Identity is delegated; authority is local.} A managed provider proves who someone is. \ProjectName{}
        decides what they may do, from the role stored in its own database; the role claim in the identity token is a
        convenience, never the source of truth.
  \item \textbf{The backend is the security boundary.} Every rule that protects people is enforced in the service
        layer; hiding a button in a client is only a courtesy.
  \item \textbf{Safety-critical rules are pure and testable.} Caller trust, identity-verification policy, webhook
        verification, role--facility policy and one-time-code handling are small classes with no database or clock
        dependency.
  \item \textbf{Side effects happen after commit.} Live updates, notifications and asynchronous work run only once the
        transaction has committed.
  \item \textbf{Fail open for emergencies, fail closed for data.} Abuse controls never refuse a report and a failing AI,
        geocoder or push provider never blocks one; conversely, a failed audit write blocks a medical-profile read.
  \item \textbf{Files never pass through the API.} Clients upload and download evidence directly to private storage with
        short-lived signed URLs.
  \item \textbf{Real time first, push second.} Open screens update over a WebSocket; push notifications reach people who
        are away from the app.
\end{itemize}

\section{Architecture}

\subsection{System context}
\ProjectName{} is a client--server system (\cref{fig:intro-context}). Two clients --- a Flutter mobile app and a Next.js web
console --- talk to one Spring Boot API over REST and a STOMP WebSocket. The API persists to PostgreSQL and integrates
with external providers for identity and push messaging (Firebase), object storage (Supabase), language-model
triage (OpenRouter), e-mail (Brevo), identity verification (Didit), and maps, routing and geocoding
(OpenStreetMap services). A separate AI model service is reached only by the backend, over an authenticated
internal interface.

\subsection{Backend structure}
The backend is a single deployable Spring Boot application organised by feature module (\cref{fig:modules}). Every
module follows the same layering: controller (HTTP and authorisation annotations), service (rules and
transactions) and repository (Spring Data JPA). Cross-cutting concerns --- token authentication, exception mapping,
scheduling and WebSocket configuration --- live beside the modules. \Cref{tab:modules} lists the responsibility of each
module and \cref{tab:jobs} the scheduled jobs that keep the system tidy.

\diagram{arch_packages.png}{Backend feature modules.}{fig:modules}
{\tabfont\rowband
\begin{longtable}{@{}>{\raggedright\arraybackslash}p{\dimexpr(\linewidth-4\tabcolsep)*211/1000\relax}>{\raggedright\arraybackslash}p{\dimexpr(\linewidth-4\tabcolsep)*789/1000\relax}@{}}
\caption{Backend feature modules and their responsibilities}\label{tab:modules}\\
\rowcolor{najdanavy}\textcolor{white}{\sffamily\bfseries\small Module} & \textcolor{white}{\sffamily\bfseries\small Responsibility}\\
\endfirsthead
\rowcolor{najdanavy}\textcolor{white}{\sffamily\bfseries\small Module} & \textcolor{white}{\sffamily\bfseries\small Responsibility}\\
\endhead
\bottomrule\endlastfoot
\textbf{auth} & E-mail OTP registration, Google/\allowbreak{}phone citizen registration, employee provisioning, e-mail availability, rate limiting.\\
\textbf{user} & Profiles, roles, address and phone rules, medical profile, responder applications, Firebase claims and reconciliation.\\
\textbf{kyc} & Didit sessions, signed webhook verification, status reconciliation, verification policy.\\
\textbf{legal} & Versioned Terms and Privacy documents, structural content validation, consent recording and the consent interceptor.\\
\textbf{security} & Firebase token filter, URL rules, bootstrap roles and super admin.\\
\textbf{incident} & Incidents, evidence, chat, missions, hospital transfers, caller trust, presence, AI triage clients, storage service, live-event publishers.\\
\textbf{unit} & Response units, shifts, facility-manager service and policies, status reconciliation.\\
\textbf{facility} & Facilities, OpenStreetMap seeding, address backfill, type-change impact, registration reconciler.\\
\textbf{hospital} & Incoming-transfer dashboard queries.\\
\textbf{notification} & Inbox, push sender, device registry, recipient resolution, bilingual messages, maintenance.\\
\textbf{ai} & HTTP client for the AI model service.\\
\textbf{guidance, content} & First-aid guidance and JSON content documents.\\
\textbf{audit, mail, status, appversion, geo, config, exceptions} & Audit log, Brevo sender, public status, release lookup, Egypt boundary and reverse geocoding, async/\allowbreak{}after-commit helpers, central exception handling.\\
\end{longtable}}

{\tabfont\rowband
\begin{longtable}{@{}>{\raggedright\arraybackslash}p{\dimexpr(\linewidth-6\tabcolsep)*263/1000\relax}>{\raggedright\arraybackslash}p{\dimexpr(\linewidth-6\tabcolsep)*184/1000\relax}>{\raggedright\arraybackslash}p{\dimexpr(\linewidth-6\tabcolsep)*553/1000\relax}@{}}
\caption{Scheduled maintenance jobs}\label{tab:jobs}\\
\rowcolor{najdanavy}\textcolor{white}{\sffamily\bfseries\small Job} & \textcolor{white}{\sffamily\bfseries\small Schedule} & \textcolor{white}{\sffamily\bfseries\small Purpose}\\
\endfirsthead
\rowcolor{najdanavy}\textcolor{white}{\sffamily\bfseries\small Job} & \textcolor{white}{\sffamily\bfseries\small Schedule} & \textcolor{white}{\sffamily\bfseries\small Purpose}\\
\endhead
\bottomrule\endlastfoot
\textbf{Email-verification cleanup} & Every 15 minutes & Deletes expired or consumed verification rows.\\
\textbf{Notification maintenance} & Daily 03:30 & Purges notifications and stale device tokens older than 90 days.\\
\textbf{KYC webhook purge} & Daily 03:30 & Deletes old webhook-event ids used for de-duplication.\\
\textbf{Facility registration reconciler} & Daily 03:45 & Recomputes the registered flag of every facility from its active staff and units.\\
\textbf{Presence pruning} & Every 15 seconds & Drops dispatcher presence entries that stopped heart-beating.\\
\textbf{Rate-limiter cleanup} & Every 5 minutes & Evicts expired in-memory rate-limit counters.\\
\end{longtable}}


\subsection{Client structure}
Both clients are layered so that screens contain no business logic (\cref{fig:clients}). On mobile, feature folders sit on
Riverpod providers and go\_router with role-aware routing, over core services for the HTTP client, the STOMP
connection, push, storage and the offline outbox. On the web, App Router routes host role workspaces built from
components, with server state cached by TanStack Query, UI state in Zustand stores and live updates arriving through
STOMP subscriptions that invalidate the relevant queries.

\diagramwide{arch_clients.png}{Layering of the Flutter mobile app and the Next.js web console.}{fig:clients}

\section{AI Components}\label{sec:ai-components}

\subsection{Report triage with a language model}
After a report is committed, two asynchronous jobs call a language model through the OpenRouter gateway with
temperature zero, so the same input gives the same output. \emph{Priority classification} turns the category, the
optional text message and the context into one of \code{CRITICAL}, \code{HIGH}, \code{MEDIUM} or \code{LOW} with a
confidence value, stored on the incident. \emph{Duplicate detection} compares the new report with open incidents of the
same category within 2\,km and 120 minutes and stores a suggested duplicate with a confidence value. Both jobs are
guarded: a time-out, a missing key or a provider error is logged and the incident is unaffected, and a dispatcher can
retry. A configuration switch disables AI entirely, and it disables itself when no key is configured.

\begin{designnote}[AI recommends, humans decide]
The model's output is displayed next to the report; it never changes what a dispatcher may do, and it never blocks
or reorders access to a report. Confirming a duplicate, in particular, is always a dispatcher's action.
\end{designnote}

\subsection{The AI model service}
A separate Python (FastAPI) service hosts versioned machine-learning pipelines behind an API key, on the internal
network only. A \emph{model} is a folder containing the fitted pipeline, a manifest that is the contract (features,
types, units, ranges, labels, metrics and hashes) and golden test cases. A registry file selects the active version and
optional shadow versions, so a model can be rolled forward or back by editing one value. At start-up each package is
verified --- the file hash against the manifest and the library versions against the training environment --- and a
faulty package is reported as degraded on the health endpoint without taking other models down.
\Cref{tab:ai-endpoints} lists its endpoints. The backend contains the matching typed client, so a trained model can be
attached to the patient-vitals stream; the first clinical model is planned work (\cref{sec:future-work}).

{\tabfont\rowband
\begin{longtable}{@{}>{\raggedright\arraybackslash}p{\dimexpr(\linewidth-6\tabcolsep)*289/1000\relax}>{\raggedright\arraybackslash}p{\dimexpr(\linewidth-6\tabcolsep)*118/1000\relax}>{\raggedright\arraybackslash}p{\dimexpr(\linewidth-6\tabcolsep)*592/1000\relax}@{}}
\caption{Endpoints of the AI model service}\label{tab:ai-endpoints}\\
\rowcolor{najdanavy}\textcolor{white}{\sffamily\bfseries\small Endpoint} & \textcolor{white}{\sffamily\bfseries\small Auth} & \textcolor{white}{\sffamily\bfseries\small Purpose}\\
\endfirsthead
\rowcolor{najdanavy}\textcolor{white}{\sffamily\bfseries\small Endpoint} & \textcolor{white}{\sffamily\bfseries\small Auth} & \textcolor{white}{\sffamily\bfseries\small Purpose}\\
\endhead
\bottomrule\endlastfoot
\textbf{GET /\allowbreak{}health} & None & Loaded models and load failures.\\
\textbf{GET /\allowbreak{}v1/\allowbreak{}models} & API key & List models and active versions.\\
\textbf{GET /\allowbreak{}v1/\allowbreak{}models/\allowbreak{}\{name\}} & API key & The manifest: exactly which features the model expects.\\
\textbf{POST /\allowbreak{}v1/\allowbreak{}models/\allowbreak{}\{name\}/\allowbreak{}predict} & API key & Validate the features (unknown names, out-of-range values and unknown categories return 422) and return label, confidence, probabilities, missing features and warnings; an optional version parameter supports shadow testing.\\
\end{longtable}}


\section{Data Model}\label{sec:persistence}

The database has 26 tables managed by Flyway migrations. Entities use \code{bigint} identity keys, most timestamps are
\code{timestamptz}, enumerations are stored as strings with check constraints, and rows that are updated
concurrently (incidents and missions) carry a \code{version} column for optimistic locking. The model is shown in two
parts that share the \code{users} entity: identity, medical and notification data (\cref{fig:erd-identity}) and the
operational model of facilities, units, incidents and missions (\cref{fig:erd-ops}). Appendix~\ref{app:database} contains the
full data dictionary.

\diagramwide{erd_identity.png}{Conceptual model of identity, medical profile, legal consent and notifications.}{fig:erd-identity}
\diagram{erd_ops.png}{Conceptual model of facilities, units, incidents, missions and hospital transfers.}{fig:erd-ops}

\subsection{Integrity rules in the schema}
Several of the system's guarantees are enforced by the database itself rather than only by application code:
\begin{itemize}
  \item a unique index on (citizen, client request id) is the database-level guarantee behind idempotent submission;
  \item a unique constraint on the personal-unit owner gives each first responder exactly one personal unit, and plate
        numbers are unique;
  \item one hospital transfer per mission, one legal acceptance per user and version, and unique version numbers per
        legal document type;
  \item indexes for the hot paths: notifications by user (newest first and unread), active units by facility, a citizen's
        incident history and audit lookups by actor and target.
\end{itemize}

\section{Class Design}

The class diagrams show the persistent domain model and the principal services. Only members that explain the design
are listed. \Cref{fig:cls-ops} presents the operational domain: users, facilities, units, shifts, incidents, evidence,
chat, missions, transfers and vitals. \Cref{fig:cls-identity} presents the identity and support domain, including the
medical profile with its embedded emergency contacts, responder applications, legal consent, notifications and audit.

\diagram{cls_ops.png}{Operational domain classes.}{fig:cls-ops}
\diagram{cls_identity.png}{Identity and support domain classes.}{fig:cls-identity}

\subsection{Dispatch services}
\code{IncidentServiceImpl} orchestrates a report's life and delegates to small collaborators: a location policy
(where a report may come from), the caller-trust service over the pure trust policy, the asynchronous AI and
geocoding services, and the mission service for cascading cancellations (\cref{fig:cls-dispatch-a}).
\code{MissionServiceImpl} takes a pessimistic lock on the unit row to serialise assignments, while optimistic version
checks protect incidents and missions; event publishers wrap every broadcast so that it runs after commit
(\cref{fig:cls-dispatch-b}).

\diagramwide{cls_dispatch_a.png}{The incident service and its collaborators.}{fig:cls-dispatch-a}
\diagram{cls_dispatch_b.png}{Mission, shift and hospital-transfer services.}{fig:cls-dispatch-b}

\subsection{Notification subsystem}
Business services call a publisher inside their transaction. The publisher resolves the recipients, builds the
bilingual text and publishes a \code{NotificationRequested} event; a dispatcher component listens \emph{after commit},
asynchronously writes the inbox rows, nudges connected clients over the WebSocket and hands push messages to a
\code{PushSender} abstraction, whose Firebase implementation prunes tokens that the provider reports as invalid
(\cref{fig:cls-notify}).

\diagram{cls_notify.png}{The notification subsystem.}{fig:cls-notify}

\subsection{Identity, verification and consent services}
Registration, one-time codes, identity verification and legal consent are separate collaborators with narrow
responsibilities (\cref{fig:cls-idsvc-a,fig:cls-idsvc-b}). The pure \code{KycPolicy} decides who must be verified;
the webhook verifier authenticates provider callbacks; and a consent interceptor protects annotated endpoints.

\diagram{cls_idsvc_a.png}{Authentication, registration and e-mail one-time-code services.}{fig:cls-idsvc-a}
\diagramwide{cls_idsvc_b.png}{Identity-verification and legal-consent services.}{fig:cls-idsvc-b}

\section{State Model}

Six entities have non-trivial life cycles. The diagrams show every transition the code allows.

\subsection{Incident}
An incident starts as \code{NEW}, may gain an AI priority (\code{AI\_PROCESSED}), becomes \code{ASSIGNED} when the first
unit is offered the work, returns to \code{DISPATCHER\_REVIEW} if its last active mission is cancelled, and ends as
\code{RESOLVED} or \code{CANCELLED} (\cref{fig:st-incident}, \cref{tab:incident-states}).

\diagramwide{st_incident.png}{Incident life cycle.}{fig:st-incident}
{\tabfont\rowband
\begin{longtable}{@{}>{\raggedright\arraybackslash}p{\dimexpr(\linewidth-6\tabcolsep)*195/1000\relax}>{\raggedright\arraybackslash}p{\dimexpr(\linewidth-6\tabcolsep)*390/1000\relax}>{\raggedright\arraybackslash}p{\dimexpr(\linewidth-6\tabcolsep)*416/1000\relax}@{}}
\caption{Incident statuses}\label{tab:incident-states}\\
\rowcolor{najdanavy}\textcolor{white}{\sffamily\bfseries\small Status} & \textcolor{white}{\sffamily\bfseries\small Meaning} & \textcolor{white}{\sffamily\bfseries\small Entered when}\\
\endfirsthead
\rowcolor{najdanavy}\textcolor{white}{\sffamily\bfseries\small Status} & \textcolor{white}{\sffamily\bfseries\small Meaning} & \textcolor{white}{\sffamily\bfseries\small Entered when}\\
\endhead
\bottomrule\endlastfoot
\textbf{NEW} & Saved, awaiting triage. & A citizen submits.\\
\textbf{AI\_\allowbreak{}PROCESSED} & Priority has been stored. & The AI classification completes.\\
\textbf{DISPATCHER\_\allowbreak{}REVIEW} & Needs a dispatcher’s decision again. & The last active mission is cancelled or withdrawn while ASSIGNED.\\
\textbf{ASSIGNED} & At least one unit has been offered work. & A dispatcher assigns the first unit.\\
\textbf{RESOLVED} & Finished. & All missions complete, or a dispatcher completes by override.\\
\textbf{CANCELLED} & Closed without a response. & Citizen cancel, confirmed duplicate, or false report before a unit arrived.\\
\end{longtable}}


\begin{caution}[A declared status that no workflow uses]
\code{IN\_PROGRESS} exists in the incident status enumeration and in the active-status sets, but no workflow enters it:
field progress is carried by the mission statuses. It is reserved for future use.
\end{caution}

\subsection{Mission and hospital transfer}
A mission is \emph{offered} to a unit and then accepted, en route and arrived, ending as completed, rejected or
cancelled. Only the unit's lead moves it forward; dispatch can cancel it at any stage. For an ambulance, the hospital
transfer runs in parallel from selection to arrival, and arrival completes the mission.

\diagramwide{st_mission.png}{Mission life cycle.}{fig:st-mission}
\diagram{st_transfer.png}{Hospital-transfer life cycle.}{fig:st-transfer}

\subsection{Response unit, identity verification and caller trust}
\diagram{st_unit.png}{Response-unit status.}{fig:st-unit}
\diagramwide{st_kyc.png}{Identity-verification status.}{fig:st-kyc}
Caller trust is derived rather than stored: each assessment recomputes a score from recent history
(\cref{sec:caller-trust}) and maps it to a tier.
\diagram{st_trust.png}{Caller-trust tiers.}{fig:st-trust}

\section{Behavioural Design}

The sequence diagrams below show the order of interactions in the most important flows.

\subsection{Account creation}
E-mail registration is deliberately a three-step process (\cref{fig:sq-register}): nothing is created in the identity
provider or the database until the address is proven, so there are no half-registered accounts to clean up. Registration
is atomic across the two stores --- if the database step fails, the new identity account is deleted again.
\diagram{sq_register.png}{E-mail registration with a one-time code.}{fig:sq-register}
\diagram{sq_google.png}{Google sign-in, profile creation and phone verification.}{fig:sq-google}
\diagram{sq_employee.png}{Administrator-provisioned employee account.}{fig:sq-employee}

\subsection{Reporting and evidence}
Submitting a report (\cref{fig:sq-submit}) does only database work before responding; AI classification, duplicate
detection and geocoding start asynchronously. Evidence travels directly between the client and storage through
signed URLs (\cref{fig:sq-evidence}), and a report taken offline is delivered later by the outbox
(\cref{fig:sq-outbox}).
\diagram{sq_submit.png}{Submitting an emergency report.}{fig:sq-submit}
\diagram{sq_evidence.png}{Uploading and viewing evidence through signed URLs.}{fig:sq-evidence}
\diagram{sq_outbox.png}{The offline outbox and idempotent delivery.}{fig:sq-outbox}

\subsection{Dispatch and field response}
Assigning a unit locks the unit, checks it, creates the mission and alerts the crew; the lead then works the mission step
by step (\cref{fig:sq-assign}). For an ambulance, the handoff to a hospital includes selecting the hospital, streaming
vitals and arriving, which closes the mission (\cref{fig:sq-hospital}).
\diagram{sq_assign.png}{Assigning a unit and working the mission.}{fig:sq-assign}
\diagram{sq_hospital.png}{Ambulance hospital handoff with live vitals.}{fig:sq-hospital}

\subsection{Notifications, identity verification and responder applications}\label{sec:after-commit}
Notifications are created and pushed only after the triggering transaction commits (\cref{fig:sq-notify}), so a
rolled-back action never produces a notification. Identity verification relies on a signed webhook that is confirmed
against the provider's own decision (\cref{fig:sq-kyc}), and responder applications end in an administrator's
approval that changes the role and creates the personal unit (\cref{fig:sq-application}).
\diagram{sq_notify.png}{Notification creation and push delivery after commit.}{fig:sq-notify}
\diagram{sq_kyc.png}{Identity verification with a signed webhook.}{fig:sq-kyc}
\diagram{sq_application.png}{First-responder application and approval.}{fig:sq-application}

\section{Authentication and Authorisation}\label{sec:authorisation}

Every request passes through the same chain: token verification, user loading, URL rules, method rules, an optional
consent check and finally ownership or scope checks inside the service (\cref{fig:sq-authz}). \Cref{tab:public-endpoints}
lists the few endpoints that do not need a token, and \cref{tab:authz-layers} summarises the authorisation layers.

\diagram{sq_authz.png}{How a request is authenticated and authorised.}{fig:sq-authz}
{\tabfont\rowband
\begin{longtable}{@{}>{\raggedright\arraybackslash}p{\dimexpr(\linewidth-4\tabcolsep)*474/1000\relax}>{\raggedright\arraybackslash}p{\dimexpr(\linewidth-4\tabcolsep)*526/1000\relax}@{}}
\caption{Endpoints that do not require a token}\label{tab:public-endpoints}\\
\rowcolor{najdanavy}\textcolor{white}{\sffamily\bfseries\small Endpoint} & \textcolor{white}{\sffamily\bfseries\small Purpose}\\
\endfirsthead
\rowcolor{najdanavy}\textcolor{white}{\sffamily\bfseries\small Endpoint} & \textcolor{white}{\sffamily\bfseries\small Purpose}\\
\endhead
\bottomrule\endlastfoot
\textbf{POST /\allowbreak{}api/\allowbreak{}auth/\allowbreak{}register/\allowbreak{}otp/\allowbreak{}send, /\allowbreak{}otp/\allowbreak{}verify, /\allowbreak{}citizen/\allowbreak{}password} & E-mail registration (rate limited).\\
\textbf{POST /\allowbreak{}api/\allowbreak{}auth/\allowbreak{}register/\allowbreak{}citizen} & Citizen profile creation for a Firebase identity (the token is read by the filter; 401 without it).\\
\textbf{GET /\allowbreak{}api/\allowbreak{}auth/\allowbreak{}email-exists} & Whether an address is registered (rate limited).\\
\textbf{POST /\allowbreak{}api/\allowbreak{}kyc/\allowbreak{}webhook} & Didit webhook; authenticated by signature, not by token.\\
\textbf{GET /\allowbreak{}api/\allowbreak{}legal/\allowbreak{}documents/\allowbreak{}\{type\}} & Current Terms and Privacy Policy.\\
\textbf{GET /\allowbreak{}api/\allowbreak{}app/\allowbreak{}version, /\allowbreak{}api/\allowbreak{}health, /\allowbreak{}api/\allowbreak{}status, /\allowbreak{}actuator/\allowbreak{}health} & Release lookup and health.\\
\textbf{/\allowbreak{}ws/\allowbreak{}**} & WebSocket handshake (origin-restricted).\\
\end{longtable}}

{\tabfont\rowband
\begin{longtable}{@{}>{\raggedright\arraybackslash}p{\dimexpr(\linewidth-6\tabcolsep)*171/1000\relax}>{\raggedright\arraybackslash}p{\dimexpr(\linewidth-6\tabcolsep)*447/1000\relax}>{\raggedright\arraybackslash}p{\dimexpr(\linewidth-6\tabcolsep)*382/1000\relax}@{}}
\caption{Authorisation layers}\label{tab:authz-layers}\\
\rowcolor{najdanavy}\textcolor{white}{\sffamily\bfseries\small Layer} & \textcolor{white}{\sffamily\bfseries\small Mechanism} & \textcolor{white}{\sffamily\bfseries\small Example}\\
\endfirsthead
\rowcolor{najdanavy}\textcolor{white}{\sffamily\bfseries\small Layer} & \textcolor{white}{\sffamily\bfseries\small Mechanism} & \textcolor{white}{\sffamily\bfseries\small Example}\\
\endhead
\bottomrule\endlastfoot
\textbf{Authentication} & Firebase token filter loads the NAJDA user; a disabled account is refused. & 401 /\allowbreak{} 403 account disabled\\
\textbf{Role} & \code{@PreAuthorize(hasRole /\allowbreak{} hasAnyRole)} on controller methods and classes. & Only dispatch roles may assign a unit.\\
\textbf{Ownership} & Service checks that the caller owns or is attached to the resource. & Only the citizen cancels their incident; only the unit’s LEAD advances a mission.\\
\textbf{Scope} & Facility scoping for managers; role--facility policies. & A manager touches only their facility’s units.\\
\textbf{Consent} & \code{@RequiresLegalConsent} interceptor on protected endpoints. & 403 LEGAL\_\allowbreak{}CONSENT\_\allowbreak{}REQUIRED until the latest terms are accepted.\\
\textbf{Data-level} & Medical-profile access service with audit. & Treating roles only, open incident only.\\
\end{longtable}}


\section{Real-Time Channels and Notifications}

The API exposes one STOMP endpoint with a simple in-memory broker and 10-second heartbeats in both directions, so an idle
connection dropped by a proxy is noticed. Most channels carry a small refresh signal and the client re-queries; a few
carry the changed object itself (\cref{tab:realtime}). Notifications (\cref{tab:notifications}) are created for each
concerned user in English and Arabic; four types are urgent and use high-priority, short-lived push on a dedicated channel.

{\tabfont\rowband
\begin{longtable}{@{}>{\raggedright\arraybackslash}p{\dimexpr(\linewidth-6\tabcolsep)*342/1000\relax}>{\raggedright\arraybackslash}p{\dimexpr(\linewidth-6\tabcolsep)*368/1000\relax}>{\raggedright\arraybackslash}p{\dimexpr(\linewidth-6\tabcolsep)*289/1000\relax}@{}}
\caption{Real-time channels}\label{tab:realtime}\\
\rowcolor{najdanavy}\textcolor{white}{\sffamily\bfseries\small Channel} & \textcolor{white}{\sffamily\bfseries\small Carries} & \textcolor{white}{\sffamily\bfseries\small Consumers}\\
\endfirsthead
\rowcolor{najdanavy}\textcolor{white}{\sffamily\bfseries\small Channel} & \textcolor{white}{\sffamily\bfseries\small Carries} & \textcolor{white}{\sffamily\bfseries\small Consumers}\\
\endhead
\bottomrule\endlastfoot
\textbf{/\allowbreak{}topic/\allowbreak{}incidents} & Incident created or updated; refresh signals. & Dispatch console, citizen screens, admin.\\
\textbf{/\allowbreak{}topic/\allowbreak{}incidents/\allowbreak{}\{id\}/\allowbreak{}chat} & A chat message. & Members of the incident.\\
\textbf{/\allowbreak{}topic/\allowbreak{}incidents/\allowbreak{}\{id\}/\allowbreak{}responders} & Responder set changed. & Citizen tracking, dispatch.\\
\textbf{/\allowbreak{}topic/\allowbreak{}incidents/\allowbreak{}\{id\}/\allowbreak{}presence} & Who is viewing the incident. & Dispatchers.\\
\textbf{/\allowbreak{}topic/\allowbreak{}missions} & A mission created or changed. & Dispatch, responders, citizen.\\
\textbf{/\allowbreak{}topic/\allowbreak{}units} & Unit positions and status. & Dispatch, responders.\\
\textbf{/\allowbreak{}topic/\allowbreak{}users} & User list changed. & Admin console.\\
\textbf{/\allowbreak{}topic/\allowbreak{}facilities} & Facility list changed. & Admin and manager consoles.\\
\textbf{/\allowbreak{}topic/\allowbreak{}hospital-transfers} & Transfers changed. & Hospital dashboard.\\
\textbf{/\allowbreak{}topic/\allowbreak{}hospital-transfers/\allowbreak{}\{id\}/\allowbreak{}vitals} & A new vitals record. & Hospital dashboard.\\
\textbf{/\allowbreak{}topic/\allowbreak{}vitals} & Vitals refresh signal. & Hospital dashboard.\\
\textbf{/\allowbreak{}topic/\allowbreak{}users/\allowbreak{}\{id\}/\allowbreak{}notifications} & Refresh signal for one user’s inbox. & That user’s clients.\\
\end{longtable}}

{\tabfont\rowband
\begin{longtable}{@{}>{\raggedright\arraybackslash}p{\dimexpr(\linewidth-8\tabcolsep)*325/1000\relax}>{\raggedright\arraybackslash}p{\dimexpr(\linewidth-8\tabcolsep)*84/1000\relax}>{\raggedright\arraybackslash}p{\dimexpr(\linewidth-8\tabcolsep)*301/1000\relax}>{\raggedright\arraybackslash}p{\dimexpr(\linewidth-8\tabcolsep)*289/1000\relax}@{}}
\caption{Notification types, recipients and triggers}\label{tab:notifications}\\
\rowcolor{najdanavy}\textcolor{white}{\sffamily\bfseries\small Type} & \textcolor{white}{\sffamily\bfseries\small Urgent} & \textcolor{white}{\sffamily\bfseries\small Recipients} & \textcolor{white}{\sffamily\bfseries\small Trigger}\\
\endfirsthead
\rowcolor{najdanavy}\textcolor{white}{\sffamily\bfseries\small Type} & \textcolor{white}{\sffamily\bfseries\small Urgent} & \textcolor{white}{\sffamily\bfseries\small Recipients} & \textcolor{white}{\sffamily\bfseries\small Trigger}\\
\endhead
\bottomrule\endlastfoot
\textbf{INCIDENT\_\allowbreak{}NEW} & Yes & All dispatchers & A report is submitted.\\
\textbf{INCIDENT\_\allowbreak{}CANCELLED\_\allowbreak{}BY\_\allowbreak{}CITIZEN} & No & Dispatchers and live crews & The citizen cancels.\\
\textbf{MISSION\_\allowbreak{}OFFERED} & Yes & The unit’s crew & A dispatcher assigns the unit.\\
\textbf{MISSION\_\allowbreak{}CANCELLED} & No & The unit’s crew & Dispatch, false report or the citizen stands the mission down.\\
\textbf{MISSION\_\allowbreak{}REJECTED} & Yes & Dispatchers & The LEAD rejects.\\
\textbf{MISSION\_\allowbreak{}WITHDRAWN} & No & Dispatchers & The LEAD withdraws.\\
\textbf{INCIDENT\_\allowbreak{}ACCEPTED /\allowbreak{} EN\_\allowbreak{}ROUTE /\allowbreak{} ARRIVED} & No & The citizen & The matching mission step.\\
\textbf{INCIDENT\_\allowbreak{}RESOLVED} & No & The citizen & The incident is resolved.\\
\textbf{INCIDENT\_\allowbreak{}CLOSED\_\allowbreak{}FALSE\_\allowbreak{}REPORT /\allowbreak{} DUPLICATE} & No & The citizen & Dispatch closes the incident.\\
\textbf{PATIENT\_\allowbreak{}INCOMING} & Yes & Staff of the destination hospital & A hospital is selected.\\
\textbf{CHAT\_\allowbreak{}MESSAGE} & No & The other side of the conversation (no message text) & A chat message is sent.\\
\textbf{APPLICATION\_\allowbreak{}SUBMITTED} & No & Admins & A responder application arrives.\\
\textbf{APPLICATION\_\allowbreak{}APPROVED /\allowbreak{} REJECTED} & No & The applicant & An Admin reviews.\\
\textbf{LEGAL\_\allowbreak{}UPDATED} & No & All users & A version requiring acceptance is published.\\
\end{longtable}}


\section{Security and Privacy Design}\label{sec:privacy}

\Cref{tab:security-design} summarises the design. Two aspects deserve elaboration.

\subsection{Privacy of the medical profile}
A patient's profile is shared only if \emph{all} of the following hold: the report shares it (never when the report is for
someone else), the patient still allows sharing, the incident is open, and the viewer is attached to the patient in a
medical capacity --- a dispatcher, the crew of an ambulance or first-responder unit on an active mission, or staff of
the hospital selected as the destination. Administrators are deliberately excluded. Every read writes an audit entry,
and the read fails if the audit entry cannot be written. Withdrawing consent takes effect immediately, even for open
incidents.

\subsection{Protecting secrets and data at rest}
Verification codes are stored only as HMAC-SHA-256 hashes and registration tokens as hashes of 256-bit random values;
comparisons are constant-time. Chat messages, text evidence, free-text medical fields, emergency-contact names and
phones, and vital-sign notes are encrypted at rest by the application, and evidence and documents live in a private
bucket reachable only by short-lived signed URLs.

{\tabfont\rowband
\begin{longtable}{@{}>{\raggedright\arraybackslash}p{\dimexpr(\linewidth-4\tabcolsep)*224/1000\relax}>{\raggedright\arraybackslash}p{\dimexpr(\linewidth-4\tabcolsep)*776/1000\relax}@{}}
\caption{Security and privacy design}\label{tab:security-design}\\
\rowcolor{najdanavy}\textcolor{white}{\sffamily\bfseries\small Concern} & \textcolor{white}{\sffamily\bfseries\small Design}\\
\endfirsthead
\rowcolor{najdanavy}\textcolor{white}{\sffamily\bfseries\small Concern} & \textcolor{white}{\sffamily\bfseries\small Design}\\
\endhead
\bottomrule\endlastfoot
\textbf{Credentials} & Passwords exist only in Firebase and are never seen by the API. Employees receive a one-time reset link; the temporary password used to create the account is never shown or stored.\\
\textbf{Registration abuse} & Per-IP and per-e-mail rate limits; resend cool-down; five-attempt limit; code and registration token stored as keyed hashes and compared in constant time.\\
\textbf{Account integrity} & Atomic cross-store registration with compensation; protected Super Admin accounts; role claim kept in step with the database role.\\
\textbf{Data at rest} & Application-layer encryption of chat, text evidence, free-text medical fields, emergency contacts and vitals notes; private object bucket.\\
\textbf{File handling} & Signed, short-lived upload and download URLs; extension allow-list; ownership and incident-state checks before issuing a URL.\\
\textbf{Medical privacy} & Standing consent plus per-report choice; server-enforced exclusion for reports on behalf of others; treating-roles-only access; admins excluded; every read audited and fail-closed.\\
\textbf{Abuse control} & Caller-trust tiers that slow but never block; false-report verdicts by humans only; manual account suspension.\\
\textbf{Webhooks} & Signature, timestamp window and event-id de-duplication; the decision is re-fetched from the provider.\\
\textbf{Consent} & Immutable versions; recorded acceptance with source; interceptor enforces latest consent.\\
\textbf{Live channel} & Origin-restricted WebSocket with heartbeats; notification text never carries chat content.\\
\textbf{Auditing} & Audit rows for medical-profile reads, false-report verdicts and clearing, phone overrides, unit delete or archive, and employee creation.\\
\end{longtable}}


\section{Behaviour When a Dependency Fails}\label{sec:failure}

Each external service is wrapped so that its failure disables only its own feature (\cref{tab:failure}). Reporting and
dispatch continue even when the language model, the geocoder, push delivery, the mail provider, the verification provider
or the object store is unavailable.

{\tabfont\rowband
\begin{longtable}{@{}>{\raggedright\arraybackslash}p{\dimexpr(\linewidth-6\tabcolsep)*237/1000\relax}>{\raggedright\arraybackslash}p{\dimexpr(\linewidth-6\tabcolsep)*197/1000\relax}>{\raggedright\arraybackslash}p{\dimexpr(\linewidth-6\tabcolsep)*566/1000\relax}@{}}
\caption{Behaviour when an external dependency fails}\label{tab:failure}\\
\rowcolor{najdanavy}\textcolor{white}{\sffamily\bfseries\small Dependency} & \textcolor{white}{\sffamily\bfseries\small If it fails} & \textcolor{white}{\sffamily\bfseries\small Effect}\\
\endfirsthead
\rowcolor{najdanavy}\textcolor{white}{\sffamily\bfseries\small Dependency} & \textcolor{white}{\sffamily\bfseries\small If it fails} & \textcolor{white}{\sffamily\bfseries\small Effect}\\
\endhead
\bottomrule\endlastfoot
\textbf{Language model (OpenRouter)} & Timeout or error & No AI priority or duplicate suggestion; incident and dispatch continue; retry available.\\
\textbf{Nominatim} & Timeout or error & Address stays empty; coordinates are still used.\\
\textbf{FCM} & Send error or invalid token & Inbox and live updates still work; invalid tokens are deleted.\\
\textbf{Brevo} & Send error & The code request fails and the visitor can try again; no account is created.\\
\textbf{Didit} & Provider error & Verification cannot start or reconcile now; reporting is unaffected unless the deployment requires verification.\\
\textbf{Supabase Storage} & Upload or signing error & Evidence cannot be attached; the report itself is unaffected.\\
\textbf{Overpass, OSRM} & Unavailable & Importing fails with a clear error (the bundled facility list is unaffected); clients fall back to straight-line distance for ETA.\\
\textbf{Firebase Authentication} & Unavailable & Sign-in and token verification fail; the public status endpoint reports authentication DOWN.\\
\end{longtable}}


\section{Deployment}

\Cref{fig:deployment} shows the deployment view. The backend runs as a container, the web console as a Next.js
application, and the AI service as an internal-only container; PostgreSQL holds the data, and the mobile app is
distributed as a signed APK. A Docker Compose file runs the whole stack locally with one command
(\cref{tab:deploy-units}); GitHub Actions builds, tests and releases (\cref{tab:pipelines}).

\diagram{deploy.png}{Deployment view.}{fig:deployment}
{\tabfont\rowband
\begin{longtable}{@{}>{\raggedright\arraybackslash}p{\dimexpr(\linewidth-6\tabcolsep)*184/1000\relax}>{\raggedright\arraybackslash}p{\dimexpr(\linewidth-6\tabcolsep)*342/1000\relax}>{\raggedright\arraybackslash}p{\dimexpr(\linewidth-6\tabcolsep)*474/1000\relax}@{}}
\caption{Deployable units and their technology}\label{tab:deploy-units}\\
\rowcolor{najdanavy}\textcolor{white}{\sffamily\bfseries\small Unit} & \textcolor{white}{\sffamily\bfseries\small Technology} & \textcolor{white}{\sffamily\bfseries\small Notes}\\
\endfirsthead
\rowcolor{najdanavy}\textcolor{white}{\sffamily\bfseries\small Unit} & \textcolor{white}{\sffamily\bfseries\small Technology} & \textcolor{white}{\sffamily\bfseries\small Notes}\\
\endhead
\bottomrule\endlastfoot
\textbf{Backend API} & Spring Boot 4, Java 21; container image (JRE, non-root user) & Runs migrations at start-up, then validates entities against the schema. Hosted as a container service (the production CORS and WebSocket origin list names the Render host and the Vercel web host).\\
\textbf{Web console} & Next.js 16, React 19; container image or Vercel & Public Firebase and API settings are provided at build time; a server route manages the session cookie.\\
\textbf{AI model service} & Python 3.12, FastAPI; container image & No published port; reachable only from the backend on the internal network.\\
\textbf{Database} & PostgreSQL 17 & Single database; Flyway baseline migration; data volume in Compose.\\
\textbf{Mobile app} & Flutter, Android; signed APK & Distributed from GitHub Releases; the API reports the newest release.\\
\textbf{Local stack} & Docker Compose & PostgreSQL, AI service, backend and web console with one command; health-checked start order (database before backend).\\
\end{longtable}}

{\tabfont\rowband
\begin{longtable}{@{}>{\raggedright\arraybackslash}p{\dimexpr(\linewidth-6\tabcolsep)*211/1000\relax}>{\raggedright\arraybackslash}p{\dimexpr(\linewidth-6\tabcolsep)*237/1000\relax}>{\raggedright\arraybackslash}p{\dimexpr(\linewidth-6\tabcolsep)*553/1000\relax}@{}}
\caption{Continuous-integration pipelines}\label{tab:pipelines}\\
\rowcolor{najdanavy}\textcolor{white}{\sffamily\bfseries\small Pipeline} & \textcolor{white}{\sffamily\bfseries\small Trigger} & \textcolor{white}{\sffamily\bfseries\small What it does}\\
\endfirsthead
\rowcolor{najdanavy}\textcolor{white}{\sffamily\bfseries\small Pipeline} & \textcolor{white}{\sffamily\bfseries\small Trigger} & \textcolor{white}{\sffamily\bfseries\small What it does}\\
\endhead
\bottomrule\endlastfoot
\textbf{AI service CI} & Every push and pull request & Installs dependencies, runs the registry and export tests, builds the Docker image.\\
\textbf{Mobile CI} & Pull request & Gets dependencies, generates localisation, runs static analysis.\\
\textbf{Build test APK} & Manual & Builds a private test APK with a test keystore after analysis.\\
\textbf{Release mobile} & Pushed version tag & Builds the production APK with the release keystore after analysis and publishes it to GitHub Releases.\\
\end{longtable}}


\section{Key Design Decisions}

The most consequential decisions are collected, with their rationale, in Appendix~\ref{app:decisions}. Four are worth stating here.

\begin{designnote}[Modular monolith, one separate AI service]
A single backend with strong internal module boundaries gives simple transactions and a simple deployment for a small
team; the AI service differs in language, lifecycle and scaling needs and is therefore deployed separately.
\end{designnote}
\begin{designnote}[Idempotent submission]
A report may be sent twice by a flaky network; the second must never create a second emergency. The client's request id
is checked first and backed by a unique index.
\end{designnote}
\begin{designnote}[Pessimistic locking for unit assignment]
Double-booking a unit is the costliest race in dispatch. A row lock serialises assignment; contention elsewhere is rare,
so optimistic version checks suffice.
\end{designnote}
\begin{designnote}[Caller trust slows but never blocks]
An emergency line that can refuse a caller can cost a life. Abuse is handled by delay, human verdicts and manual
suspension.
\end{designnote}

\section{Chapter Summary}

This chapter described the architecture, the AI components, the data model, the classes, the state models, the
behaviour of the main flows, the authentication and authorisation model, real-time messaging and notifications, the
security and privacy design, failure behaviour and deployment. \Cref{chap:implementation} explains how these designs were
realised in code.
